\section{Haag Kastler Axioms}\label{sctn:haag-kastler-axioms}
With that out of the way we can state the axioms. Each axiom below is presented as a \emph{definition} so that it is captured as a node in the blueprint declaration graph. In the Lean formalization each axiom corresponds to a \texttt{Prop}-valued predicate (or, in the case of \emph{Local Algebras}, a data field), and the bundling structure \texttt{HaagKastlerNet} packages them together (see \ref{def:haag-kastler-net}). Downstream theorems take an instance of \texttt{HaagKastlerNet} as a hypothesis and invoke each axiom as a projection.

\begin{definition}[Axiom 1: Local Algebras]
    \label{def:local-algebras}
    \lean{Physicslib4.AQFT.HaagKastler.LocalNet}
    \leanfile{Physicslib4/AQFT/HaagKastler/LocalAlgebras.lean}
    \leanok
    \uses{def:alexandrov-topology, def:minkowski-spacetime, def:chronological-future-and-chronological-past}
    For any basis element $\mathbf{B}$ of the Alexandrov topology on Minkowski spacetime, i.e. any set of the form $I^+(p) \cap I^-(q)$, there is a corresponding abstract C*-algebra $\mathfrak{U}(\mathbf{B})$
    \begin{align}
        \mathbf{B} \longmapsto \mathfrak{U}(\mathbf{B}),
    \end{align}
    and when $\mathbf{B}$ is the empty set, we have the distinguished correspondence
    \begin{align}
        \emptyset \longmapsto \mathbb{C} \mathbf{1},
    \end{align}
    where $\mathbf{1}$ is the multiplicative identity in the abstract C*-algebra $\mathbb{C} \mathbf{1}$.
\end{definition}

\begin{definition}[Axiom 2: Isotony]
    \label{def:isotony}
    \lean{Physicslib4.AQFT.HaagKastler.Isotony}
    \leanfile{Physicslib4/AQFT/HaagKastler/Isotony.lean}
    \leanok
    \uses{def:alexandrov-topology, def:minkowski-spacetime, def:chronological-future-and-chronological-past, def:local-algebras}
    Let $\mathbf{B}_1$ and $\mathbf{B}_2$ be any two basis elements of the Alexandrov topology on Minkowski spacetime, i.e. any two sets of the form $I^+(p_1) \cap I^-(q_1)$ and $I^+(p_2) \cap I^-(q_2)$.

    If $\mathbf{B}_1 \subset \mathbf{B}_2$ then $\mathfrak{U}(\mathbf{B}_1) \subset \mathfrak{U}(\mathbf{B}_2)$, where inclusion is implemented by
    \begin{align}
        i : \mathfrak{U}(\mathbf{B}_1) \hookrightarrow \mathfrak{U}(\mathbf{B}_2),
    \end{align}
    a unital *-monomorphism.
\end{definition}

Before introducing the next axiom, we must introduce the definition:

\begin{definition}[Quasilocal Algebra]
    \label{def:quasilocal-algebra}
    \lean{Physicslib4.AQFT.HaagKastler.QuasilocalAlgebra}
    \leanfile{Physicslib4/AQFT/HaagKastler/QuasilocalAlgebra.lean}
    \leanok
    \uses{def:local-algebras}
    Consider the set-theoretic union of all $\mathfrak{U}(\mathbf{B})$. As previously proven, this set-theoretic union is a normed *-algebra. Also, as previously proven, taking its completion one obtains a C*-algebra denoted as $\mathfrak{U}$. This C*-algebra $\mathfrak{U}$ is called the \textit{quasilocal algebra}.
\end{definition}

\begin{definition}[Axiom 3: Local Commutativity]
    \label{def:local-commutativity}
    \lean{Physicslib4.AQFT.HaagKastler.LocalCommutativity}
    \leanfile{Physicslib4/AQFT/HaagKastler/LocalCommutativity.lean}
    \leanok
    \uses{def:alexandrov-topology, def:minkowski-spacetime, def:chronological-future-and-chronological-past, def:completely-spacelike, def:local-algebras, def:quasilocal-algebra}
    Let $\mathbf{B}_1$ and $\mathbf{B}_2$ be any two basis elements of the Alexandrov topology on Minkowski spacetime, i.e. any two sets of the form $I^+(p_1) \cap I^-(q_1)$ and $I^+(p_2) \cap I^-(q_2)$.

    If $\mathbf{B_1}$ and $\mathbf{B_2}$ are completely spacelike, then $\mathfrak{U}(\mathbf{B_1})$ and $\mathfrak{U}(\mathbf{B_2})$ commute in the quasilocal algebra $\mathfrak{U}$, i.e. for any $a_1$ in $\mathfrak{U}(\mathbf{B_1})$ and $a_2$ in $\mathfrak{U}(\mathbf{B_2})$ it follows that
    \begin{align}
        i(a_1) i(a_2) - i(a_2) i(a_1) = 0
    \end{align}
    in the quasilocal algebra $\mathfrak{U}$.
\end{definition}

The next axiom makes use of the following new definition:

\begin{definition}[Quasilocal Observable]
    \label{def:quasilocal-observable}
    \lean{Physicslib4.AQFT.HaagKastler.IsQuasilocalObservable}
    \leanfile{Physicslib4/AQFT/HaagKastler/QuasilocalCompleteness.lean}
    \leanok
    \uses{def:quasilocal-algebra, thrm:gns-construction-theorem, def:state}
    The image $\pi_\omega(a)$ of a self-adjoint member $a$ of the quasilocal algebra $\mathfrak{U}$ under a GNS *-homomorphism $\pi_\omega$ is self-adjoint and thus corresponds to an ``observable''. Any ``observable'' corresponding to such a self-adjoint $\pi_\omega(a)$ is called a \textit{quasilocal observable}.
\end{definition}

\begin{definition}[Axiom 4: Quasilocal Completeness]
    \label{def:quasilocal-completeness}
    \lean{Physicslib4.AQFT.HaagKastler.QuasilocalCompleteness}
    \leanfile{Physicslib4/AQFT/HaagKastler/QuasilocalCompleteness.lean}
    \leanok
    \uses{def:quasilocal-observable}
    All ``observables'' are quasilocal observables.
\end{definition}

\begin{definition}[Axiom 5: Lorentz Covariance]
    \label{def:lorentz-covariance}
    \lean{Physicslib4.AQFT.HaagKastler.LorentzCovariance}
    \leanfile{Physicslib4/AQFT/HaagKastler/LorentzCovariance.lean}
    \leanok
    \uses{def:alexandrov-topology, def:minkowski-spacetime, def:chronological-future-and-chronological-past, def:local-algebras, def:isotony}
    Let $\mathbf{B}$ be any basis element of the Alexandrov topology on Minkowski spacetime, i.e. any set of the form $I^+(p) \cap I^-(q)$.

    A member $L$ of the inhomogeneous Lorentz group connected to the identity acts on $\mathfrak{U}(\mathbf{B})$ as follows
    \begin{align}
        \alpha_L : \mathfrak{U}(\mathbf{B}) &\longrightarrow \mathfrak{U}(L\mathbf{B}) \\
                             a              &\longmapsto     \alpha_L(a)
    \end{align}
    where $L\mathbf{B}$ is the image of the region $\mathbf{B}$ under the transformation $L$ and $\alpha_L$ is a unital *-isomorphism generated by $L$. The map $\alpha_L$ is such that (1) for the identity element $\mathbf{1}$ of the Lorentz group it satisfies
    \begin{align}
        \alpha_{\mathbf{1}}(a) = a,
    \end{align}
    (2) for all appropriate $a$, $L$, and $L'$ it satisfies
    \begin{align}
        \alpha_{L' \cdot L}(a) = \alpha_{L'}(\alpha_{L}(a)),
    \end{align}
    and (3) for basis elements $\mathbf{B}_\iota \subset \mathbf{B}_\kappa$ and the unital *-monomorphism $i$ of Axiom 2 (Isotony) $\alpha_L$ commutes with $i$. In other words the following diagram

    \begin{center}
        \begin{tikzcd}
            \mathfrak{U}(\mathbf{B}_\kappa) \arrow[r, "\alpha_L"]
            & \mathfrak{U}(L\mathbf{B}_\kappa) \\
            \mathfrak{U}(\mathbf{B}_\iota) \arrow[r, "\alpha_L"]  \arrow[u, "i"]
            & \mathfrak{U}(L\mathbf{B}_\iota)  \arrow[u, "i"]
        \end{tikzcd}
    \end{center}

    commutes.
\end{definition}

\begin{definition}[Haag-Kastler Net]
    \label{def:haag-kastler-net}
    \lean{Physicslib4.AQFT.HaagKastler.HaagKastlerNet}
    \leanfile{Physicslib4/AQFT/HaagKastler/Net.lean}
    \leanok
    \uses{def:local-algebras, def:isotony, def:local-commutativity, def:quasilocal-completeness, def:lorentz-covariance}
    A \emph{Haag-Kastler net} on Minkowski spacetime is the bundling of the data of \ref{def:local-algebras} together with the properties of \ref{def:isotony}, \ref{def:local-commutativity}, \ref{def:quasilocal-completeness}, and \ref{def:lorentz-covariance}. In the Lean formalization this is a single structure whose fields are the assignment $\mathbf{B} \mapsto \mathfrak{U}(\mathbf{B})$ and proofs that this assignment satisfies the four remaining axioms. Theorems about AQFT take an instance of this structure as a hypothesis and invoke each axiom as a projection.
\end{definition}

This concludes the presentation of the ``sharpened'' axioms. As proven in this blog post, these ``sharpened'' axioms entail the original axioms save Axiom 6 (Primitivity), which is an axiom that we abandon.

These ``sharpened'' axioms also allow for a straightforward generalization to a set of axioms describing AQFT in curved spacetime, i.e. a set of axioms describing AQFT in a curved spacetime background that is fixed and treated classically. In a subsequent blog post we will detail these axioms.

\subsection{Einstein Causality}

Local commutativity (\ref{def:local-commutativity}) is an algebraic statement about the quasilocal algebra. Its physical content - that spacelike-separated measurements do not interfere - is its operator form: in any representation of the quasilocal algebra, spacelike-separated local observables commute as bounded operators on the Hilbert space.

\begin{theorem}[Einstein Causality in a Representation]
    \label{thrm:einstein-causality}
    \lean{Physicslib4.AQFT.HaagKastler.HaagKastlerNet.einstein_causality, Physicslib4.AQFT.HaagKastler.HaagKastlerNet.exists_gns_einstein_causality}
    \leanfile{Physicslib4/AQFT/HaagKastler/EinsteinCausality.lean}
    \leanok
    \uses{def:local-commutativity, def:haag-kastler-net, thrm:gns-construction-theorem, def:state}
    Let $\pi$ be any $*$-representation of the quasilocal algebra on a Hilbert space $H$. If $\mathbf{B}_1, \mathbf{B}_2$ are completely spacelike-separated basis regions, then for all $a \in \mathfrak{U}(\mathbf{B}_1)$ and $b \in \mathfrak{U}(\mathbf{B}_2)$ the operators $\pi(\iota_{\mathbf{B}_1} a)$ and $\pi(\iota_{\mathbf{B}_2} b)$ commute. This is the image under $\pi$ of the local commutativity relation (\ref{def:local-commutativity}); since the GNS representation of any state is such a $\pi$, spacelike-separated local observables commute on every GNS Hilbert space.
\end{theorem}

\subsection{Local von Neumann Algebras}

The abstract C*-algebras $\mathfrak{U}(\mathbf{B})$ are a stepping stone; the genuine object of algebraic quantum field theory is the net of von Neumann algebras they generate in a representation. Given a $*$-representation $\pi$ of the quasilocal algebra on a Hilbert space $H$, the \emph{local von Neumann algebra} of a region $\mathbf{B}$ is the bicommutant $R(\mathbf{B}) = \pi(\mathfrak{U}(\mathbf{B}))''$ of the local observable operators. We model the commutant by the centralizer in $\mathcal{B}(H)$.

\begin{definition}[Local von Neumann Algebra]
    \label{def:local-von-neumann}
    \lean{Physicslib4.AQFT.HaagKastler.HaagKastlerNet.localOperators, Physicslib4.AQFT.HaagKastler.HaagKastlerNet.localVonNeumann}
    \leanfile{Physicslib4/AQFT/HaagKastler/LocalVonNeumann.lean}
    \leanok
    \uses{def:haag-kastler-net, thrm:gns-construction-theorem}
    Let $\pi$ be a $*$-representation of the quasilocal algebra on $H$. The \emph{local observable operators} of a region $\mathbf{B}$ are the image $\pi(\mathfrak{U}(\mathbf{B})) = \{\pi(\iota_{\mathbf{B}} a) : a \in \mathfrak{U}(\mathbf{B})\}$. The \emph{local von Neumann algebra} $R(\mathbf{B})$ is the bicommutant $\pi(\mathfrak{U}(\mathbf{B}))''$.
\end{definition}

\begin{definition}[$R(\mathbf{B})$ as a von Neumann Algebra]
    \label{def:local-von-neumann-algebra}
    \lean{Physicslib4.AQFT.HaagKastler.HaagKastlerNet.localVonNeumannAlgebra, Physicslib4.GNS.vonNeumannOfSelfAdjoint}
    \leanfile{Physicslib4/AQFT/HaagKastler/LocalVonNeumann.lean}
    \leanok
    \uses{def:local-von-neumann}
    The local algebra $R(\mathbf{B})$ is registered as a genuine \texttt{VonNeumannAlgebra} (Mathlib's bundled structure), not merely a set of operators. The general fact is that the bicommutant $S''$ of any self-adjoint set $S$ of bounded operators is a von Neumann algebra: $S'$ is a $*$-subalgebra (the centralizer of a self-adjoint set is $*$-closed), and $S'' = S''''$ since the triple commutant collapses. The local observable operators $\pi(\mathfrak{U}(\mathbf{B}))$ are self-adjoint ($\pi$ and $\iota$ are $*$-homomorphisms), so $R(\mathbf{B}) = \pi(\mathfrak{U}(\mathbf{B}))''$ is a von Neumann algebra whose underlying set is the bicommutant of \ref{def:local-von-neumann}.
\end{definition}

\begin{theorem}[Microcausality at the von Neumann Level]
    \label{thrm:von-neumann-microcausality}
    \lean{Physicslib4.AQFT.HaagKastler.HaagKastlerNet.localVonNeumann_subset_centralizer}
    \leanfile{Physicslib4/AQFT/HaagKastler/LocalVonNeumann.lean}
    \leanok
    \uses{def:local-von-neumann, thrm:einstein-causality}
    For completely spacelike-separated basis regions $\mathbf{B}_1, \mathbf{B}_2$, the local von Neumann algebras commute: $R(\mathbf{B}_1) \subseteq R(\mathbf{B}_2)'$. This is the von Neumann form of Einstein causality (\ref{thrm:einstein-causality}): elementwise commutation $\pi(\mathfrak{U}(\mathbf{B}_2)) \subseteq \pi(\mathfrak{U}(\mathbf{B}_1))'$, with the antitonicity of the commutant and the identity $S''' = S'$ collapsing the iterated commutants. A Haag-Kastler net is thus a net of mutually commuting von Neumann algebras for spacelike regions.
\end{theorem}

\begin{theorem}[Isotony of the von Neumann Net]
    \label{thrm:von-neumann-isotony}
    \lean{Physicslib4.AQFT.HaagKastler.HaagKastlerNet.localVonNeumann_mono}
    \leanfile{Physicslib4/AQFT/HaagKastler/LocalVonNeumann.lean}
    \leanok
    \uses{def:local-von-neumann, def:isotony}
    For basis regions $\mathbf{B}_1 \subseteq \mathbf{B}_2$, the local von Neumann algebras are nested: $R(\mathbf{B}_1) \subseteq R(\mathbf{B}_2)$. The local observables of $\mathbf{B}_1$ embed into those of $\mathbf{B}_2$ via the quasilocal isotony coherence, and the double commutant is monotone.
\end{theorem}

\begin{theorem}[Bundled von Neumann Microcausality and Isotony]
    \label{thrm:von-neumann-bundled-order}
    \lean{Physicslib4.AQFT.HaagKastler.HaagKastlerNet.localVonNeumannAlgebra_le_commutant, Physicslib4.AQFT.HaagKastler.HaagKastlerNet.localVonNeumannAlgebra_mono}
    \leanfile{Physicslib4/AQFT/HaagKastler/LocalVonNeumann.lean}
    \leanok
    \uses{def:local-von-neumann-algebra, thrm:von-neumann-microcausality, thrm:von-neumann-isotony}
    Phrased on the bundled \texttt{VonNeumannAlgebra} $R(\mathbf{B})$ with its order $\le$ and Mathlib's commutant: for completely spacelike-separated regions $R(\mathbf{B}_1) \le R(\mathbf{B}_2)'$, and for $\mathbf{B}_1 \subseteq \mathbf{B}_2$, $R(\mathbf{B}_1) \le R(\mathbf{B}_2)$. Both reduce to the set-level statements through the coercion $\uparrow R(\mathbf{B}) = \pi(\mathfrak{U}(\mathbf{B}))''$.
\end{theorem}

\begin{theorem}[Statistical Independence (Schlieder Property)]
    \label{thrm:statistical-independence}
    \lean{Physicslib4.AQFT.HaagKastler.HaagKastlerNet.localVonNeumann_separating, Physicslib4.AQFT.HaagKastler.HaagKastlerNet.eq_zero_of_commute_of_cyclic}
    \leanfile{Physicslib4/AQFT/HaagKastler/LocalVonNeumann.lean}
    \leanok
    \uses{thrm:von-neumann-microcausality, def:cyclic-vector}
    If $\Omega$ is cyclic for the local observables of $\mathbf{B}_1$, then for a spacelike-separated region $\mathbf{B}_2$ every $R \in R(\mathbf{B}_2)$ with $R\Omega = 0$ is zero, so $\Omega$ is separating for $R(\mathbf{B}_2)$. The implication is elementary ($R$ vanishes on the dense set of vectors $A\Omega$ for $A$ a local observable of $\mathbf{B}_1$, since $R(\mathbf{B}_2)$ commutes with those observables by microcausality, \ref{thrm:von-neumann-microcausality}). In Minkowski spacetime the cyclicity hypothesis is exactly the content of the Reeh-Schlieder theorem, which rests on the spectrum condition; here it is taken as an explicit hypothesis, mirroring the curved-spacetime version where no spectrum condition is available.
\end{theorem}

\begin{theorem}[Statistical Independence, bundled]
    \label{thrm:statistical-independence-bundled}
    \lean{Physicslib4.AQFT.HaagKastler.HaagKastlerNet.localVonNeumannAlgebra_separating}
    \leanfile{Physicslib4/AQFT/HaagKastler/LocalVonNeumann.lean}
    \leanok
    \uses{thrm:statistical-independence, def:local-von-neumann-algebra}
    The separating property phrased on the bundled \texttt{VonNeumannAlgebra} $R(\mathbf{B}_2)$: with $\Omega$ cyclic for the local observables of $\mathbf{B}_1$, every $R$ in the bundled algebra $R(\mathbf{B}_2)$ of a spacelike-separated region with $R\Omega = 0$ is zero. It reduces to \ref{thrm:statistical-independence} through the coercion $\uparrow R(\mathbf{B}_2) = \pi(\mathfrak{U}(\mathbf{B}_2))''$.
\end{theorem}

\begin{theorem}[Geometric Covariance of the von Neumann Net]
    \label{thrm:von-neumann-geometric-covariance}
    \lean{Physicslib4.AQFT.HaagKastler.CovariantQuasilocalAlgebra.lieConj_image_covLocalVonNeumann, Physicslib4.AQFT.HaagKastler.CovariantQuasilocalAlgebra.lieConj_image_covLocalOperators}
    \leanfile{Physicslib4/AQFT/HaagKastler/GeometricCovariance.lean}
    \leanok
    \uses{def:local-von-neumann, def:covariant-quasilocal-algebra, thrm:irreducible-covariant-representation}
    In a covariant representation $\pi$ of the quasilocal algebra with the operator covariance $U(L)\,\pi(a)\,U(L)^{-1} = \pi(\beta_L a)$, conjugation by the implementing unitary $U(L)$ carries the local von Neumann algebra of a region $\mathbf{B}$ onto that of the Lorentz-transformed region $L \cdot \mathbf{B}$:
    \[ U(L)\, R(\mathbf{B})\, U(L)^{-1} = R(L \cdot \mathbf{B}). \]
    This is the statement that the symmetry group acts \emph{geometrically} on the net of von Neumann algebras. The proof rests on three ingredients, all already available: operator covariance (the last clause of \ref{thrm:irreducible-covariant-representation}); the fact that $\beta_L = $ the covariance action sends $\iota_{\mathbf{B}}(\mathfrak{U}(\mathbf{B}))$ onto $\iota_{L\cdot\mathbf{B}}(\mathfrak{U}(L\cdot\mathbf{B}))$ (from the action-on-generators identity and surjectivity of the covariance equivalence $\alpha_L$); and the algebraic fact that conjugation by a unit is a multiplicative automorphism, which therefore commutes with the bicommutant. The reusable core is that a multiplicative automorphism maps the centralizer of a set onto the centralizer of its image, hence maps bicommutants to bicommutants. In particular $R(\mathbf{B})$ and $R(L \cdot \mathbf{B})$ are unitarily equivalent (conjugate by $U(L)$), so any unitary-conjugation-invariant property of a local von Neumann algebra --- such as being a factor --- is constant along the Lorentz orbit of a region.
\end{theorem}

\begin{theorem}[Orbit-Invariance of Factoriality]
    \label{thrm:von-neumann-factor-orbit}
    \lean{Physicslib4.AQFT.HaagKastler.CovariantQuasilocalAlgebra.covLocalVonNeumann_isFactor_smul}
    \leanfile{Physicslib4/AQFT/HaagKastler/GeometricCovariance.lean}
    \leanok
    \uses{thrm:von-neumann-geometric-covariance, thrm:irreducible-factor}
    If the local von Neumann algebra $R(\mathbf{B})$ is a \emph{factor} (its center $R(\mathbf{B}) \cap R(\mathbf{B})'$ is exactly the scalars), then so is $R(L \cdot \mathbf{B})$ for every Lorentz transformation $L$. Geometric covariance (\ref{thrm:von-neumann-geometric-covariance}) exhibits $R(L \cdot \mathbf{B}) = U(L) R(\mathbf{B}) U(L)^{-1}$; conjugation by a unitary is a multiplicative automorphism that carries the center onto the center and fixes the scalars, so it preserves the factor property. Thus being a factor is constant along the Lorentz orbit of a region.
\end{theorem}

\subsection{Irreducibility and Schur's Lemma}

A representation is \emph{irreducible} when its commutant is trivial: the only operators commuting with every $\pi(a)$ are scalars. This is the von Neumann (commutant) form of irreducibility. The cornerstone, connecting the commutant to the GNS state, is the topological Schur lemma for a cyclic representation.

\begin{definition}[Irreducible Representation]
    \label{def:irreducible-representation}
    \lean{Physicslib4.GNS.IsIrreducible}
    \leanfile{Physicslib4/GNS/Irreducibility.lean}
    \leanok
    A $*$-representation $\pi : A \to \mathcal{B}(H)$ is \emph{irreducible} if every bounded operator $T$ commuting with all $\pi(a)$ is a scalar multiple of the identity, $T = c \cdot 1$.
\end{definition}

\begin{theorem}[Topological Schur Lemma]
    \label{thrm:schur-lemma}
    \lean{Physicslib4.GNS.eq_smul_one_of_commute_of_cyclic}
    \leanfile{Physicslib4/GNS/Irreducibility.lean}
    \leanok
    \uses{def:cyclic-vector}
    Let $\Omega$ be a cyclic vector for $\pi$. If $T$ commutes with every $\pi(a)$ and the diagonal coefficient $a \mapsto \langle \Omega, T\,\pi(a)\Omega\rangle$ equals $c$ times $a \mapsto \langle \Omega, \pi(a)\Omega\rangle$, then $T = c \cdot 1$. The proof is pure Hilbert-space analysis: using the $*$-representation property and the commutation relation, every off-diagonal coefficient $\langle \pi(b)\Omega, (T - c)\,\pi(a)\Omega\rangle$ vanishes, so $(T - c)$ annihilates the dense cyclic orbit and is therefore zero.
\end{theorem}

\begin{theorem}[Commutant Scalar iff Proportional Coefficient]
    \label{thrm:commutant-scalar-iff}
    \lean{Physicslib4.GNS.isScalar_iff_coeff_proportional}
    \leanfile{Physicslib4/GNS/Irreducibility.lean}
    \leanok
    \uses{thrm:schur-lemma, def:state, thrm:gns-construction-theorem}
    In a cyclic representation reproducing a state $\omega$, an operator $T$ commuting with all $\pi(a)$ is a scalar multiple of the identity if and only if its diagonal coefficient $a \mapsto \langle \Omega, T\,\pi(a)\Omega\rangle$ is a scalar multiple of $\omega$. This is the precise operator-theoretic bridge to purity: irreducibility (every commutant element is scalar) is exactly the statement that every commutant coefficient is proportional to $\omega$.
\end{theorem}

\begin{definition}[Pure State]
    \label{def:pure-state}
    \lean{Physicslib4.GNS.IsPure}
    \leanfile{Physicslib4/GNS/Irreducibility.lean}
    \leanok
    \uses{def:state}
    A state $\omega$ is \emph{pure} if every positive linear functional $\psi$ dominated by $\omega$ (that is, $0 \le \psi(a^*a) \le \omega(a^*a)$ for all $a$) is a scalar multiple of $\omega$. This order-theoretic characterization is the extreme-point notion of purity, phrased to avoid convex-combination and normalization bookkeeping.
\end{definition}

One direction of the classical equivalence ``$\omega$ pure $\iff$ GNS representation irreducible'' is established below.

\begin{theorem}[Pure Implies Irreducible]
    \label{thrm:pure-implies-irreducible}
    \lean{Physicslib4.GNS.isIrreducible_of_isPure}
    \leanfile{Physicslib4/GNS/Irreducibility.lean}
    \leanok
    \uses{def:pure-state, def:irreducible-representation, thrm:commutant-scalar-iff}
    If $\omega$ is pure, then any cyclic representation reproducing $\omega$ (in particular its GNS representation) is irreducible. Since the commutant is $*$-closed, a commuting operator decomposes into self-adjoint real and imaginary parts; for a self-adjoint commuting $S$, the affine rescaling $T = (2(\Vert S\Vert+1))^{-1} S + \tfrac12$ satisfies $0 \le T \le 1$, so its coefficient functional $a \mapsto \langle \Omega, T\,\pi(a)\Omega\rangle$ is positive and dominated by $\omega$. Purity forces it proportional to $\omega$, whence $T$ (and so $S$, and so the original operator) is a scalar by \ref{thrm:commutant-scalar-iff}.
\end{theorem}

The converse, irreducible $\Rightarrow$ pure, requires the other half of the GNS Radon-Nikodym correspondence: that every dominated positive functional $\psi \le \omega$ is represented by an operator $T$ in the commutant, via the bounded sesquilinear form $(\pi(a)\Omega, \pi(b)\Omega) \mapsto \psi(a^* b)$. The analytic crux is that this form is well-defined and bounded.

\begin{theorem}[The GNS Radon-Nikodym Form is Bounded]
    \label{thrm:gns-form-bound}
    \lean{Physicslib4.GNS.gns_form_norm_le, Physicslib4.GNS.gns_form_well_defined}
    \leanfile{Physicslib4/GNS/RadonNikodym.lean}
    \leanok
    \uses{def:pure-state, lmm:cauchy-schwarz-inequality}
    For a positive functional $\psi$ dominated by the state $\omega$ in its cyclic GNS representation, the form values obey $\Vert\psi(a^* b)\Vert \le \Vert\pi(a)\Omega\Vert\,\Vert\pi(b)\Omega\Vert$. This follows from the Cauchy-Schwarz inequality for $\psi$, the domination $\psi \le \omega$, and the reproducing identity $\omega(x^* x) = \Vert\pi(x)\Omega\Vert^2$. In particular the form depends only on the GNS vectors $\pi(a)\Omega, \pi(b)\Omega$, not on the representatives $a, b$, so it is well-defined on the cyclic subspace.
\end{theorem}

\begin{theorem}[The GNS Radon-Nikodym Operator]
    \label{thrm:gns-radon-nikodym-operator}
    \lean{Physicslib4.GNS.rnOp, Physicslib4.GNS.rnOp_inner, Physicslib4.GNS.rnOp_commute, Physicslib4.GNS.rnOp_reproducing}
    \leanfile{Physicslib4/GNS/RadonNikodym.lean}
    \leanok
    \uses{thrm:gns-form-bound}
    The bounded form of \ref{thrm:gns-form-bound} is the inner product of an operator $T$ on the GNS space: extending the densely-defined bounded form by the Fréchet-Riesz representation and a norm-controlled dense extension yields $T$ with $\langle \pi(a)\Omega, T\,\pi(b)\Omega\rangle = \psi(a^* b)$. From this reproducing identity, $T$ commutes with every $\pi(c)$ (an adjoint computation on the dense cyclic vectors), and $\psi(a) = \langle \Omega, T\,\pi(a)\Omega\rangle$ (taking $a = 1$ in the first slot).
\end{theorem}

\begin{theorem}[Pure $\iff$ Irreducible]
    \label{thrm:pure-iff-irreducible}
    \lean{Physicslib4.GNS.isPure_iff_isIrreducible}
    \leanfile{Physicslib4/GNS/RadonNikodym.lean}
    \leanok
    \uses{thrm:pure-implies-irreducible, thrm:gns-radon-nikodym-operator, thrm:commutant-scalar-iff}
    A state $\omega$ is pure if and only if its cyclic GNS representation is irreducible. For the converse direction, given a dominated $\psi$, its Radon-Nikodym operator $T$ commutes with $\pi$, so by irreducibility $T = c \cdot 1$; the reproducing identity then gives $\psi(a) = c\,\omega(a)$, so $\psi$ is a scalar multiple of $\omega$ and $\omega$ is pure.
\end{theorem}

Irreducibility has a von Neumann algebraic reading: the representation generates a \emph{factor}.

\begin{theorem}[An Irreducible Representation Generates a Factor]
    \label{thrm:irreducible-factor}
    \lean{Physicslib4.GNS.center_gnsVonNeumann_eq_of_isIrreducible}
    \leanfile{Physicslib4/GNS/Irreducibility.lean}
    \leanok
    \uses{def:irreducible-representation, def:local-von-neumann}
    The von Neumann algebra $\pi(A)''$ generated by an irreducible representation has \emph{trivial center}: an operator lies in the center $\pi(A)'' \cap (\pi(A)'')'$ if and only if it is a scalar multiple of the identity. Indeed the center is contained in $(\pi(A)'')' = \pi(A)'$ (the triple commutant collapses to the single one), which irreducibility makes the scalars; conversely scalars are central.
\end{theorem}

\begin{theorem}[The GNS Representation of a Pure State is a Factor]
    \label{thrm:pure-factor}
    \lean{Physicslib4.GNS.exists_gns_factor_of_isPure}
    \leanfile{Physicslib4/GNS/RadonNikodym.lean}
    \leanok
    \uses{def:pure-state, thrm:pure-iff-irreducible, thrm:irreducible-factor, thrm:gns-construction-theorem}
    For a pure state $\omega$ there is a cyclic GNS triple $(H, \pi, \Omega)$ reproducing $\omega$ whose generated von Neumann algebra $\pi(A)''$ has trivial center, i.e. is a factor. This combines the GNS construction, purity $\Rightarrow$ irreducibility (\ref{thrm:pure-iff-irreducible}), and \ref{thrm:irreducible-factor}. It applies verbatim to the quasilocal algebra $\mathfrak{U}$ and to each local algebra $\mathfrak{U}(\mathbf{B})$, since both are C*-algebras.
\end{theorem}

The order-theoretic definition of purity (\ref{def:pure-state}) coincides with the convex-geometric one: $\omega$ is an extreme point of the state space. The analytic key is the normalization identity for positive functionals.

\begin{theorem}[Norm of a Positive Functional]
    \label{thrm:norm-positive-functional}
    \lean{Physicslib4.GNS.norm_eq_re_apply_one_of_positive}
    \leanfile{Physicslib4/GNS/ExtremeState.lean}
    \leanok
    \uses{lmm:cauchy-schwarz-inequality}
    For a positive linear functional $\varphi$ on a unital C*-algebra, $\Vert\varphi\Vert = \mathrm{Re}\,\varphi(1)$. The bound $\mathrm{Re}\,\varphi(1) \le \Vert\varphi\Vert$ is immediate from $\Vert 1\Vert = 1$; conversely Cauchy-Schwarz with the first slot equal to $1$ gives $\Vert\varphi(b)\Vert^2 \le \mathrm{Re}\,\varphi(1)\cdot\mathrm{Re}\,\varphi(b^* b) \le \mathrm{Re}\,\varphi(1)\cdot\Vert\varphi\Vert\,\Vert b\Vert^2$, whence $\Vert\varphi\Vert^2 \le \mathrm{Re}\,\varphi(1)\cdot\Vert\varphi\Vert$. In particular every state satisfies $\omega(1) = 1$.
\end{theorem}

\begin{definition}[Extreme Point of the State Space]
    \label{def:extreme-state}
    \lean{Physicslib4.GNS.State.IsExtremePoint}
    \leanfile{Physicslib4/GNS/ExtremeState.lean}
    \leanok
    \uses{def:state}
    A state $\omega$ is an \emph{extreme point} of the (convex) state space if it is not a nontrivial convex combination of two distinct states: whenever $\omega = t\,\omega_1 + (1-t)\,\omega_2$ with $0 < t < 1$ and $\omega_1, \omega_2$ states, then $\omega_1 = \omega_2$.
\end{definition}

\begin{theorem}[Pure $\iff$ Extreme Point]
    \label{thrm:pure-iff-extreme}
    \lean{Physicslib4.GNS.isPure_iff_isExtremePoint}
    \leanfile{Physicslib4/GNS/ExtremeState.lean}
    \leanok
    \uses{def:pure-state, def:extreme-state, thrm:norm-positive-functional}
    A state $\omega$ is pure if and only if it is an extreme point of the state space. If $\omega$ is pure and $\omega = t\,\omega_1 + (1-t)\,\omega_2$, then $t\,\omega_1$ is a positive functional dominated by $\omega$, hence (by purity) a scalar multiple of $\omega$; evaluating at $1$, where every state gives $1$, pins the scalar to $t$ and forces $\omega_1 = \omega = \omega_2$. Conversely, if $\omega$ is extreme and $\psi \le \omega$ is dominated, set $\lambda = \mathrm{Re}\,\psi(1) \in [0,1]$; for $\lambda \in (0,1)$ the rescaled functionals $\lambda^{-1}\psi$ and $(1-\lambda)^{-1}(\omega - \psi)$ are states (normalized via \ref{thrm:norm-positive-functional}) whose convex combination is $\omega$, so extremality identifies them and forces $\psi = \lambda\,\omega$; the boundary cases $\lambda = 0$ and $\lambda = 1$ give $\psi = 0$ and $\psi = \omega$ respectively.
\end{theorem}

\begin{theorem}[State-Space Convexity and the Extreme-Point Bridge]
    \label{thrm:state-space-convex-bridge}
    \lean{Physicslib4.GNS.convex_stateSpace, Physicslib4.GNS.mem_extremePoints_iff_isExtremePoint}
    \leanfile{Physicslib4/GNS/ExtremeState.lean}
    \leanok
    \uses{def:state, def:extreme-state}
    Realizing the state space as the subset $\mathrm{stateSpace}(A) \subseteq A \to_L \mathbb{C}$ (the latter a real topological vector space via $\texttt{NormedSpace.complexToReal}$), it is convex: a real convex combination $a\,\omega_1 + b\,\omega_2$ of states is again a state, via $\texttt{State.convexCombo}$. Moreover a state $\omega$ lies in Mathlib's $\mathrm{extremePoints}_{\mathbb{R}}$ of the state space exactly when it is extreme in the sense of \ref{def:extreme-state}, connecting the purity characterizations to the Krein-Milman/Choquet API.
\end{theorem}

These equivalences specialize to the quasilocal algebra $\mathfrak{U}$ of a Minkowski net, where they characterize purity of a global state - the natural setting for the vacuum and other distinguished states.

\begin{theorem}[Pure $\iff$ Extreme Point on the Quasilocal Algebra]
    \label{thrm:pure-iff-extreme-quasilocal}
    \lean{Physicslib4.AQFT.HaagKastler.HaagKastlerNet.pure_iff_extreme}
    \leanfile{Physicslib4/AQFT/HaagKastler/Purity.lean}
    \leanok
    \uses{def:pure-state, def:extreme-state, thrm:pure-iff-extreme, def:quasilocal-completeness}
    A state $\omega$ on the canonical quasilocal algebra $\mathfrak{U}$ of a Haag-Kastler net is pure if and only if it is an extreme point of the state space of $\mathfrak{U}$. This is the abstract equivalence \ref{thrm:pure-iff-extreme} applied to the C*-algebra $\mathfrak{U}$.
\end{theorem}

\begin{theorem}[Pure $\iff$ Irreducible GNS on the Quasilocal Algebra]
    \label{thrm:pure-iff-irreducible-quasilocal}
    \lean{Physicslib4.AQFT.HaagKastler.HaagKastlerNet.exists_gns_pure_iff_irreducible}
    \leanfile{Physicslib4/AQFT/HaagKastler/Purity.lean}
    \leanok
    \uses{def:pure-state, def:irreducible-representation, thrm:pure-iff-irreducible, thrm:gns-construction-theorem, def:quasilocal-completeness}
    For a state $\omega$ on the quasilocal algebra $\mathfrak{U}$ there is a GNS triple $(H, \pi, \Omega)$ reproducing $\omega$ in which $\omega$ is pure if and only if the representation $\pi$ is irreducible. This combines the GNS construction with the abstract \ref{thrm:pure-iff-irreducible}.
\end{theorem}

\subsection{Covariant States and the Covariance Action}

The Lorentz covariance of a Haag-Kastler net (\ref{def:lorentz-covariance}) acts on the net fiberwise, through the $*$-isomorphisms $\alpha_L : \mathfrak{U}(\mathbf{B}) \to \mathfrak{U}(L\cdot\mathbf{B})$. We record two consequences: the notion of a Lorentz-covariant family of local states, and the lift of the fiberwise action to a single dynamical $*$-automorphism of the quasilocal algebra.

\begin{definition}[Covariant Family of Local States]
    \label{def:covariant-state-family}
    \lean{Physicslib4.AQFT.HaagKastler.HaagKastlerNet.IsCovariantFamily}
    \leanfile{Physicslib4/AQFT/HaagKastler/CovariantState.lean}
    \leanok
    \uses{def:haag-kastler-net, def:state, def:lorentz-covariance}
    Given a Haag-Kastler net (\ref{def:haag-kastler-net}), a \emph{covariant family of local states} assigns to every region $\mathbf{B}$ a state $\omega_{\mathbf{B}}$ on the local algebra $\mathfrak{U}(\mathbf{B})$ such that, for every Lorentz transformation $L$ and every $a \in \mathfrak{U}(\mathbf{B})$, one has $\omega_{\mathbf{B}}(a) = \omega_{L\cdot\mathbf{B}}(\alpha_L a)$, where $\alpha_L$ is the covariance isomorphism of \ref{def:lorentz-covariance}. The local states are thus intertwined by the Lorentz action.
\end{definition}

\begin{lemma}[Composition of Covariance]
    \label{lmm:covariant-state-family-compose}
    \lean{Physicslib4.AQFT.HaagKastler.HaagKastlerNet.IsCovariantFamily.comp}
    \leanfile{Physicslib4/AQFT/HaagKastler/CovariantState.lean}
    \leanok
    \uses{def:covariant-state-family}
    For a covariant family of local states, the covariance relation composes along the group: $\omega_{\mathbf{B}}(a) = \omega_{L'\cdot(L\cdot\mathbf{B})}\big(\alpha_{L'}(\alpha_L a)\big)$, reflecting the multiplicativity of the Lorentz action.
\end{lemma}

\begin{definition}[Quasilocal Covariance Automorphism]
    \label{def:quasilocal-lift}
    \lean{Physicslib4.AQFT.HaagKastler.HaagKastlerNet.QuasilocalLift}
    \leanfile{Physicslib4/AQFT/HaagKastler/QuasilocalAction.lean}
    \leanok
    \uses{def:quasilocal-algebra, def:lorentz-covariance}
    A \emph{lift} of the fiberwise Lorentz action of $L$ to a quasilocal algebra $\mathfrak{U}$ (\ref{def:quasilocal-algebra}) is a $*$-automorphism $\beta_L$ of $\mathfrak{U}$ intertwining the local embeddings $\iota_{\mathbf{B}}$ with the covariance isomorphisms: $\beta_L(\iota_{\mathbf{B}} a) = \iota_{L\cdot\mathbf{B}}(\alpha_L a)$ for every Alexandrov-basis set $\mathbf{B}$.
\end{definition}

\begin{lemma}[Uniqueness of the Quasilocal Lift]
    \label{lmm:quasilocal-lift-unique}
    \lean{Physicslib4.AQFT.HaagKastler.HaagKastlerNet.QuasilocalLift.unique}
    \leanfile{Physicslib4/AQFT/HaagKastler/QuasilocalAction.lean}
    \leanok
    \uses{def:quasilocal-lift}
    Any two lifts of the same $L$ have the same underlying automorphism: they agree on the union of the local images, which is dense in $\mathfrak{U}$, and $*$-automorphisms of a C*-algebra are continuous.
\end{lemma}

\begin{theorem}[Existence of the Quasilocal Lift]
    \label{thrm:quasilocal-lift-exists}
    \lean{Physicslib4.AQFT.HaagKastler.nonempty_quasilocalLift}
    \leanfile{Physicslib4/AQFT/HaagKastler/QuasilocalIntertwiner.lean}
    \leanok
    \uses{def:quasilocal-lift, def:quasilocal-algebra}
    For a covariance-compatible quasilocal algebra, the fiberwise Lorentz action extends to a $*$-automorphism of $\mathfrak{U}$, so the lift of \ref{def:quasilocal-lift} exists. The intertwiner is defined on the directed union of local images (a dense $*$-subalgebra) and extended by uniform continuity; the inverse is supplied by $L^{-1}$, and the group laws $\beta_{L'L} = \beta_{L'}\circ\beta_L$ and $\beta_{1} = \mathrm{id}$ furnish the two-sided inverse.
\end{theorem}

\begin{theorem}[Existence for the Trivial Net]
    \label{thrm:quasilocal-lift-trivial}
    \lean{Physicslib4.AQFT.HaagKastler.nonempty_trivialQuasilocalLift}
    \leanfile{Physicslib4/AQFT/HaagKastler/QuasilocalIntertwiner.lean}
    \leanok
    \uses{thrm:quasilocal-lift-exists}
    The covariance-compatibility hypothesis is satisfiable: the trivial net's quasilocal algebra is covariance-compatible (every $*$-automorphism of $\mathbb{C}$ is the identity), so the quasilocal lift exists unconditionally for the trivial net.
\end{theorem}

\begin{definition}[Covariant Quasilocal Algebra]
    \label{def:covariant-quasilocal-algebra}
    \lean{Physicslib4.AQFT.HaagKastler.CovariantQuasilocalAlgebra}
    \leanfile{Physicslib4/AQFT/HaagKastler/QuasilocalIntertwiner.lean}
    \leanok
    \uses{def:haag-kastler-net, def:quasilocal-algebra, def:lorentz-covariance, def:quasilocal-lift, thrm:quasilocal-lift-exists}
    A \emph{covariant quasilocal algebra} bundles a Haag-Kastler net (\ref{def:haag-kastler-net}), a quasilocal algebra of that net (\ref{def:quasilocal-algebra}), and a proof that the embeddings are covariance-compatible for every Lorentz transformation. On such data the quasilocal lift (\ref{def:quasilocal-lift}) exists for every $L$ by \ref{thrm:quasilocal-lift-exists}, yielding the covariance action $L \mapsto \beta_L$ on the quasilocal algebra; the trivial net provides an instance. This is the natural home for the covariance dynamics: the compatibility hypothesis of \ref{thrm:quasilocal-lift-exists} becomes structural data rather than a side condition.
\end{definition}

\begin{lemma}[Group-Action Coherence of the Covariance Automorphism]
    \label{lmm:quasilocal-action-coherence}
    \lean{Physicslib4.AQFT.HaagKastler.CovariantQuasilocalAlgebra.action_one, Physicslib4.AQFT.HaagKastler.CovariantQuasilocalAlgebra.action_mul}
    \leanfile{Physicslib4/AQFT/HaagKastler/QuasilocalIntertwiner.lean}
    \leanok
    \uses{def:covariant-quasilocal-algebra}
    The covariance action $L \mapsto \beta_L$ of a covariant quasilocal algebra is a genuine action of the Lorentz group by $*$-automorphisms of the quasilocal algebra: $\beta_{\mathbf{1}} = \mathrm{id}$ and $\beta_{L'L} = \beta_{L'} \circ \beta_L$.
\end{lemma}

\begin{definition}[Invariant State]
    \label{def:invariant-state}
    \lean{Physicslib4.AQFT.HaagKastler.CovariantQuasilocalAlgebra.IsInvariantState}
    \leanfile{Physicslib4/AQFT/HaagKastler/QuasilocalIntertwiner.lean}
    \leanok
    \uses{def:covariant-quasilocal-algebra, def:state}
    A state $\omega$ on the quasilocal algebra of a covariant quasilocal algebra (\ref{def:covariant-quasilocal-algebra}) is \emph{(Poincar\'e-)invariant} if it is a fixed point of the dual covariance action: $\omega(\beta_L a) = \omega(a)$ for every Lorentz transformation $L$ and observable $a$. This is the invariance condition of a vacuum state; further conditions (e.g.\ the spectrum condition) are imposed separately.
\end{definition}

\begin{theorem}[GNS Unitary Implementation of an Invariant State]
    \label{thrm:invariant-state-gns-unitary}
    \lean{Physicslib4.AQFT.HaagKastler.CovariantQuasilocalAlgebra.IsInvariantState.exists_gns_unitary}
    \leanfile{Physicslib4/AQFT/HaagKastler/QuasilocalIntertwiner.lean}
    \leanok
    \uses{def:invariant-state, thrm:gns-construction-theorem, lmm:quasilocal-action-coherence}
    For an invariant state $\omega$, the covariance action is implemented on the GNS Hilbert space by a family of unitaries $U(L)$ satisfying $U(L)\,\pi(a)\Omega = \pi(\beta_L a)\,\Omega$ and $U(L)\Omega = \Omega$. The unitaries are obtained by extending the densely-defined isometry $\pi(a)\Omega \mapsto \pi(\beta_L a)\Omega$ - isometric because $\omega$ preserves the GNS inner product $\langle \pi(a)\Omega, \pi(b)\Omega\rangle = \omega(a^* b)$ - to the whole GNS space.
\end{theorem}

Combining invariance with purity makes the GNS representation both covariant and irreducible. We emphasise that this is a \emph{precursor} to a vacuum representation, not a vacuum itself: a genuine vacuum additionally requires the spectrum condition (positivity of the energy-momentum spectrum), which is not imposed - and indeed not expressible - here.

\begin{theorem}[Irreducible Covariant Representation of a Pure Invariant State]
    \label{thrm:irreducible-covariant-representation}
    \lean{Physicslib4.AQFT.HaagKastler.CovariantQuasilocalAlgebra.IsInvariantState.exists_gns_irreducible_covariant}
    \leanfile{Physicslib4/AQFT/HaagKastler/QuasilocalIntertwiner.lean}
    \leanok
    \uses{thrm:invariant-state-gns-unitary, def:pure-state, def:irreducible-representation, thrm:pure-iff-irreducible}
    A state $\omega$ on the quasilocal algebra that is both invariant under the covariance action and pure yields a GNS representation that is simultaneously \emph{covariant} - implemented by a unitary representation $U(L)$ of the inhomogeneous Lorentz group fixing the cyclic vector $\Omega$, with operator covariance $U(L)\,\pi(a)\,U(L)^{-1} = \pi(\beta_L a)$ - and \emph{irreducible}. This combines the covariant GNS triple of \ref{thrm:invariant-state-gns-unitary} with purity $\Rightarrow$ irreducibility (\ref{thrm:pure-iff-irreducible}). It is a necessary precursor to a vacuum representation; the spectrum condition would be the remaining ingredient.
\end{theorem}

\begin{theorem}[Purity is Covariance-Invariant]
    \label{thrm:purity-covariance-invariant}
    \lean{Physicslib4.GNS.isPure_precomp_iff, Physicslib4.AQFT.HaagKastler.CovariantQuasilocalAlgebra.isPure_precomp_action_iff}
    \leanfile{Physicslib4/GNS/ExtremeState.lean}
    \leanok
    \uses{def:pure-state, lmm:quasilocal-action-coherence}
    Purity is preserved by any $*$-automorphism: for $\Phi : \mathfrak{U} \xrightarrow{\sim} \mathfrak{U}$, the pullback state $\omega \circ \Phi$ is pure if and only if $\omega$ is. A dominated positive functional $\psi \le \omega \circ \Phi$ transports to $\psi \circ \Phi^{-1} \le \omega$, which purity sends to a scalar multiple of $\omega$; transporting back gives $\psi$ proportional to $\omega \circ \Phi$. Applied to the covariance automorphism $\Phi = \beta_L$, this says purity of a state is a Lorentz-covariance-invariant property.
\end{theorem}

\subsection{The Separating Vector of a Faithful State}

A faithful state has a sharper consequence at the level of the GNS construction than the faithfulness of its representation: its cyclic vector is also \emph{separating}. This is the basic datum of the modular (Tomita-Takesaki) theory of the associated von Neumann algebra.

\begin{theorem}[Separating Vector of a Faithful State]
    \label{thrm:separating-faithful}
    \lean{Physicslib4.GNS.separating_of_faithful, Physicslib4.GNS.exists_gns_separating}
    \leanfile{Physicslib4/GNS/Separating.lean}
    \leanok
    \uses{def:state, def:cyclic-vector, thrm:gns-construction-theorem}
    Let $\omega$ be a faithful state and $(\mathcal{H}_\omega, \pi_\omega, \Omega)$ its GNS triple. Then the cyclic vector $\Omega$ is \emph{separating} for $\pi_\omega(\mathfrak{U})$: if $\pi_\omega(a)\Omega = 0$ then $a = 0$. Indeed $\pi_\omega(a)\Omega = 0$ gives $\omega(a^* a) = \langle \Omega, \pi_\omega(a^* a)\Omega\rangle = \langle \pi_\omega(a)\Omega, \pi_\omega(a)\Omega\rangle = 0$, and faithfulness forces $a = 0$. This holds in any representation reproducing a faithful state, not only the GNS one.
\end{theorem}

\subsection{The KMS Condition and Thermal Equilibrium}

Poincar\'e (or Killing-flow) invariance is only part of what singles out a physical equilibrium state. The Kubo-Martin-Schwinger (KMS) condition is the algebraic characterization of thermal equilibrium at inverse temperature $\beta$. It is phrased purely as an \emph{analyticity} statement about correlation functions, so - unlike the spectrum condition - it needs no unbounded-operator theory (no Stone theorem, no spectral measures). It is exactly the condition satisfied by the curved-spacetime examples of \ref{thrm:gns-unitary-stabilizer-in-curved-spacetime}: the Hartle-Hawking state on a black-hole exterior and the Gibbons-Hawking state in the de Sitter static patch are KMS for the relevant Killing flow.

\begin{definition}[One-Parameter Automorphism Group]
    \label{def:one-parameter-aut}
    \lean{Physicslib4.AQFT.IsOneParameterAut}
    \leanfile{Physicslib4/AQFT/KMS.lean}
    \leanok
    A family $\alpha : \mathbb{R} \to (A \simeq_{\star\mathrm{a}} A)$ of $*$-automorphisms of a C$^*$-algebra $A$ is a \emph{one-parameter group} if $\alpha_0 = \mathrm{id}$ and $\alpha_{s+t} = \alpha_s \circ \alpha_t$. This is the algebraic time evolution; in the curved-spacetime setting it is the automorphism group induced by a Killing flow.
\end{definition}

\begin{definition}[KMS State]
    \label{def:kms-state}
    \lean{Physicslib4.AQFT.IsKMSState}
    \leanfile{Physicslib4/AQFT/KMS.lean}
    \leanok
    \uses{def:one-parameter-aut, def:state}
    A state $\omega$ on $A$ is \emph{$(\alpha, \beta)$-KMS} for a one-parameter automorphism group $\alpha$ (\ref{def:one-parameter-aut}) at inverse temperature $\beta$ if for every $a, b \in A$ the correlation function $t \mapsto \omega(a\,\alpha_t b)$ extends to a function $F$ on the closed strip $0 \le \operatorname{Im} z \le \beta$ that is continuous there, holomorphic on the open strip, bounded, and whose boundary value on $\operatorname{Im} z = \beta$ is $t \mapsto \omega(\alpha_t b\, \cdot a)$.
\end{definition}

\begin{theorem}[The KMS State Set is Convex]
    \label{thrm:kms-convex}
    \lean{Physicslib4.AQFT.IsKMSState.convexCombo}
    \leanfile{Physicslib4/AQFT/KMS.lean}
    \leanok
    \uses{def:kms-state}
    A convex combination $s\,\omega_1 + (1-s)\,\omega_2$ ($0 \le s \le 1$) of two $(\alpha, \beta)$-KMS states is again $(\alpha, \beta)$-KMS: for each pair $(a, b)$ the analytic interpolant is the convex combination $s\,F_1 + (1-s)\,F_2$ of the two interpolants, which is continuous, bounded, and holomorphic on the strip, with boundary values that add. Thus the equilibrium states at a fixed temperature form a convex set.
\end{theorem}

\begin{lemma}[Boundary Coincidence for $a = 1$]
    \label{lmm:kms-correlation-one}
    \lean{Physicslib4.AQFT.IsKMSState.correlationOne}
    \leanfile{Physicslib4/AQFT/KMS.lean}
    \leanok
    \uses{def:kms-state}
    For a KMS state $\omega$ and any observable $a$, the correlation function $F$ of the pair $(1, a)$ has its two boundary values \emph{equal} - both are $t \mapsto \omega(\alpha_t a)$. This is the algebraic heart of the invariance argument: it follows directly from the KMS condition with $a := 1$, since $\omega(1\cdot\alpha_t a) = \omega(\alpha_t a \cdot 1) = \omega(\alpha_t a)$.
\end{lemma}

\begin{definition}[Strip-Liouville Principle]
    \label{def:strip-liouville}
    \lean{Physicslib4.AQFT.StripLiouville}
    \leanfile{Physicslib4/AQFT/KMS.lean}
    \leanok
    The \emph{strip-Liouville principle} at width $\beta$ is the statement that any function $F$ continuous and bounded on the closed strip $0 \le \operatorname{Im} z \le \beta$, holomorphic on the open strip, and with equal boundary values $F(t) = F(t + i\beta)$ for all real $t$, is constant along the real axis: $F(t) = F(0)$. It is the analytic input that turns boundary coincidence into invariance.
\end{definition}

\begin{theorem}[$i\beta$-Periodic Entire Extension (Strip Schwarz Reflection)]
    \label{thrm:strip-periodic-extension}
    \lean{Physicslib4.exists_bounded_entire_extension_of_strip_periodic}
    \leanfile{Physicslib4/Analysis/StripPeriodicExtension.lean}
    \leanok
    A function $F$ continuous on the closed strip $0 \le \operatorname{Im} z \le \beta$, holomorphic on the open strip, bounded, and with equal boundary values $F(t) = F(t + i\beta)$ on the real axis, admits a \emph{bounded entire} extension $H$ agreeing with $F$ on $\mathbb{R}$. The extension is the $i\beta$-periodic continuation $H(z) = F\!\left(z - \lfloor \operatorname{Im} z / \beta \rfloor\, i\beta\right)$, which folds every point into the fundamental strip: it is continuous across each gluing line $\operatorname{Im} z = k\beta$ because the boundary values match by periodicity, holomorphic off the lines as $F$ composed with a holomorphic shift, and holomorphic on the lines by the horizontal-line removable-singularity theorem (Morera). This is the analytic engine behind the strip-Liouville principle.
\end{theorem}

\begin{theorem}[Strip-Liouville Holds for $\beta > 0$]
    \label{thrm:strip-liouville-pos}
    \lean{Physicslib4.AQFT.stripLiouville_of_pos}
    \leanfile{Physicslib4/AQFT/KMS.lean}
    \leanok
    \uses{def:strip-liouville, thrm:strip-periodic-extension}
    At positive width $\beta > 0$ the strip-Liouville principle (\ref{def:strip-liouville}) is a theorem. The equal boundary values let $F$ extend to a bounded $i\beta$-periodic \emph{entire} function by the strip Schwarz reflection (\ref{thrm:strip-periodic-extension}), which is then constant on $\mathbb{R}$ by Liouville's theorem. The hypothesis $\beta > 0$ is necessary: at $\beta = 0$ the open strip is empty and at $\beta < 0$ the strip itself is empty, and in both cases the principle is false.
\end{theorem}

\begin{theorem}[KMS States are Invariant]
    \label{thrm:kms-invariance}
    \lean{Physicslib4.AQFT.IsKMSState.invariant_of_pos}
    \leanfile{Physicslib4/AQFT/KMS.lean}
    \leanok
    \uses{def:kms-state, lmm:kms-correlation-one, thrm:strip-liouville-pos}
    At positive inverse temperature $\beta > 0$, every $(\alpha, \beta)$-KMS state $\omega$ is $\alpha$-invariant: $\omega(\alpha_t a) = \omega(a)$ for all $t$ and $a$. By boundary coincidence (\ref{lmm:kms-correlation-one}) the KMS correlation function of $(1, a)$ has equal boundary values $F(t) = F(t + i\beta) = \omega(\alpha_t a)$; the strip-Liouville principle for $\beta > 0$ (\ref{thrm:strip-liouville-pos}) forces $F(t) = F(0)$, so $\omega(\alpha_t a) = F(t) = F(0) = \omega(\alpha_0 a) = \omega(a)$.
\end{theorem}

\begin{theorem}[Uniqueness on the Strip from Boundary Values]
    \label{thrm:strip-uniqueness}
    \lean{Physicslib4.eqOn_strip_of_eq_boundary}
    \leanfile{Physicslib4/Analysis/StripPeriodicExtension.lean}
    \leanok
    \uses{thrm:strip-periodic-extension}
    Two functions continuous and bounded on the closed strip $0 \le \operatorname{Im} z \le \beta$ (with $\beta > 0$), holomorphic on the open strip, that agree on \emph{both} boundary lines $\operatorname{Im} z = 0$ and $\operatorname{Im} z = \beta$, agree on the whole strip. Indeed their difference has vanishing boundary values, so its $i\beta$-periodic extension (\ref{thrm:strip-periodic-extension}) is a bounded entire function, constant by Liouville, equal to its value $0$ at the origin.
\end{theorem}

\begin{theorem}[Uniqueness of the KMS Correlation Function]
    \label{thrm:kms-correlation-unique}
    \lean{Physicslib4.AQFT.IsKMSState.correlation_eqOn}
    \leanfile{Physicslib4/AQFT/KMS.lean}
    \leanok
    \uses{def:kms-state, thrm:strip-uniqueness}
    For $\beta > 0$, the analytic completion of a KMS correlation function is unique: any two functions satisfying the KMS analytic data for the same pair $(a, b)$ - continuous and bounded on the strip, holomorphic on the interior, with the prescribed boundary values $t \mapsto \omega(a\,\alpha_t b)$ and $t \mapsto \omega(\alpha_t b\,\cdot a)$ - agree on the whole strip. This is the strip uniqueness (\ref{thrm:strip-uniqueness}) applied to the two analytic completions, which share both boundary values.
\end{theorem}

\subsection{KMS States for the Covariance Flow}

A one-parameter subgroup $t \mapsto L_t$ of the inhomogeneous Lorentz group - for instance the time-translation subgroup - induces, via the quasilocal lift $\beta_L$ (\ref{lmm:quasilocal-action-coherence}), a one-parameter group of $*$-automorphisms of the global quasilocal algebra $\mathfrak{U}$. Unlike the curved case, where the absence of a global algebra forces a restriction to the stabilizer $\mathrm{Stab}(\mathbf{B})$, here $\beta_L$ is a genuine automorphism of the single algebra $\mathfrak{U}$ for every $L$, so no restriction is needed.

\begin{definition}[Covariance-Flow Automorphism Family]
    \label{def:flow-aut-covariance}
    \lean{Physicslib4.AQFT.HaagKastler.CovariantQuasilocalAlgebra.flowAut}
    \leanfile{Physicslib4/AQFT/HaagKastler/QuasilocalKMS.lean}
    \leanok
    \uses{def:covariant-quasilocal-algebra, lmm:quasilocal-action-coherence}
    Given a one-parameter subgroup $t \mapsto L_t$ of the inhomogeneous Lorentz group, the \emph{covariance-flow automorphism family} of $\mathfrak{U}$ is $t \mapsto \beta_{L_t}$, the covariance action evaluated along the flow.
\end{definition}

\begin{lemma}[A Lorentz One-Parameter Subgroup Induces a One-Parameter Group]
    \label{lmm:one-parameter-aut-flow-covariance}
    \lean{Physicslib4.AQFT.HaagKastler.CovariantQuasilocalAlgebra.isOneParameterAut_flowAut}
    \leanfile{Physicslib4/AQFT/HaagKastler/QuasilocalKMS.lean}
    \leanok
    \uses{def:flow-aut-covariance, def:one-parameter-aut}
    If $t \mapsto L_t$ is a one-parameter subgroup ($L_0 = 1$, $L_{s+t} = L_s L_t$), then the induced automorphisms $t \mapsto \beta_{L_t}$ of $\mathfrak{U}$ form a one-parameter automorphism group (\ref{def:one-parameter-aut}), by the coherence of the covariance action (\ref{lmm:quasilocal-action-coherence}).
\end{lemma}

\begin{definition}[KMS State for the Covariance Flow]
    \label{def:kms-state-for-flow-covariance}
    \lean{Physicslib4.AQFT.HaagKastler.CovariantQuasilocalAlgebra.IsKMSStateForFlow}
    \leanfile{Physicslib4/AQFT/HaagKastler/QuasilocalKMS.lean}
    \leanok
    \uses{lmm:one-parameter-aut-flow-covariance, def:kms-state, def:state}
    A state $\omega$ on the quasilocal algebra $\mathfrak{U}$ is a \emph{KMS state for the covariance flow} $L$ at inverse temperature $\beta$ if it satisfies the KMS condition (\ref{def:kms-state}) for the induced one-parameter automorphism group (\ref{lmm:one-parameter-aut-flow-covariance}).
\end{definition}

\begin{theorem}[Convexity of the Covariance-Flow KMS States]
    \label{thrm:kms-convex-for-flow-covariance}
    \lean{Physicslib4.AQFT.HaagKastler.CovariantQuasilocalAlgebra.IsKMSStateForFlow.convexCombo}
    \leanfile{Physicslib4/AQFT/HaagKastler/QuasilocalKMS.lean}
    \leanok
    \uses{def:kms-state-for-flow-covariance, thrm:kms-convex}
    A convex combination $s\,\omega_1 + (1-s)\,\omega_2$ ($0 \le s \le 1$) of two KMS states on $\mathfrak{U}$ for the same covariance flow $L$ at the same inverse temperature $\beta$ is again a KMS state for that flow. This specializes the abstract KMS convexity (\ref{thrm:kms-convex}) to the induced one-parameter group $\mathrm{flowAut}$. Physically, the equilibrium states for a one-parameter symmetry flow form a convex set.
\end{theorem}
